\section{Model and irreducible geometric variables}

\label{sec:2}

\subsection{Coframe and action}

Take spatial indices $i,j,k=0,1,2$, circle index $3$ and orientation $\epsilon^{0123}=+1$. The circle coordinate has period one. The left-invariant forms obey $d\sigma^i=2\epsilon_{ijk}\sigma^j\wedge\sigma^k$, with the spatial pair summed. For $r,L>0$ and a real matrix $M$ of determinant one, define

\begin{equation}\tag{M1}\label{eq:M1}
\begin{aligned}
e^i&=rM^i{}_j\sigma^j,& e^3&=L(dy+\omega_i\sigma^i),\\
e'^a&=O^a{}_b e^b,\quad O\in SO(4),&
g&=\delta_{ab}e'^a e'^b,\qquad
\frac{d\mathcal V_4}{dy}=\frac{\pi^2Lr^3}{4}.
\end{aligned}
\end{equation}

The spatially integrated four-volume density in (M1) integrates to $\pi^2Lr^3/4$ for constant $r,L$. The physical round radius is $r/2$. The three mixing parameters are $\delta_i=\tan\chi_i$ in $O=O_{03}(\chi_0)O_{13}(\chi_1)O_{23}(\chi_2)$. They rotate the coframe without changing the metric. Geometric torsion is specified in the resulting frame by

\begin{equation}\tag{M2}\label{eq:M2}
\begin{aligned}
T_{ijk}&=\frac{2\eta}{r}\epsilon_{ijk},\\
T^{(V)}_{abc}&=\frac13(\delta_{ac}V_b-\delta_{ab}V_c),
\qquad V_b=V\delta_{b3},\\
\mathcal S&=-\int\mathrm{vol}_4
\left(\frac{R}{2\kappa^2}+\theta N\right),\qquad
N\,\mathrm{vol}_4=d(e'^a\wedge T_a).
\end{aligned}
\end{equation}

The axial term is zero if any index is 3. AX means $V=0$, VT means $\eta=0$, and MX allows both. The coupling $\theta$ is constant, and no curvature-squared term is included. The connection is metric compatible and determined by this coframe and geometric torsion. Appendix A states its components; Appendix I gives the sign map to the code accompanying Papers01/02. Prescribing torsion components in the rotated frame is distinct from simultaneously rotating all fields as a gauge transformation.

Our sign convention extends the reduced-potential convention $U=-S$ of Papers01/02: on the common EC plus complete-NY model, $\mathcal S=-S_{\mathrm{01,FULL}}=-S_{\mathrm{02}}|_{\alpha=0}$. Thus $\mathcal S=U$ on constant profiles. All lower-bound statements below concern $U$ or $\mathcal S$ with this sign.

\subsection{The common static potential}

For constant $M,\omega,O,r,L,\eta,V$, set $S=MM^{\mathsf T}$. Cartan contraction gives the static action $U$, with $F$ collecting its shape dependence:

\begin{equation}\tag{M3}\label{eq:M3}
\begin{aligned}
U&=\mathcal K\left[F(S)+\frac{3\eta^2}{4}
+\frac{V^2r^2}{12}+\frac{L^2}{r^2}\lVert M\omega\rVert^2\right],\\
F(S)&=2\mathrm{tr}(S^2)-[\mathrm{tr}(S)]^2,\qquad
\mathcal K=\frac{\pi^2Lr}{\kappa^2}>0.
\end{aligned}
\end{equation}

Complete NY vanishes pointwise on this constant family, as proved in Section 6. The scalar action is independent of the mixing rotation: the torsion norms are invariant and their divergence vanishes. The coframe is global on $S^3\times S^1$, so for smooth periodic circle variations the integral of complete NY also vanishes, although its local density can be nonzero.

\subsection{Representation and variation families}

The scalar variables are $(r,L,\eta,V)$. Twist $\omega_i$ and the three mixing coordinates transform as spatial vectors near the round frame. A volume-preserving shape can be parametrized by $M=e^h$, with $h$ symmetric and traceless; its five normalized components form spin-2. At an axial reference shape we use $M=M_0e^h$ instead. Appendix C lists the basis. At the round point, rotation invariance separates different irreducible representations in the static quadratic expansion; scalar variables may still mix with one another.

Static extrema concern (M3) with designated coordinates varied at fixed controls. For the circle functional all coordinates may be smooth periodic functions of $y$, with $r,L>0$ and $\det M=1$. The scalar restriction sets $M=I,\omega=0$; the vector restriction fixes $r,L,\eta,V$ and sets $M=I$. Both follow from the joint action below. We keep their fixed data explicit when stating stationary or minimizing properties.

\subsection{The joint circle action}

Let $J=Lr^3$. The induced three-metric, lapse and shift in the invariant spatial basis are

\begin{equation}\tag{M4}\label{eq:M4}
\gamma=r^2M^{\mathsf T}M+L^2\omega\omega^{\mathsf T},
\qquad \mathcal N^2=\frac{J^2}{\det\gamma},
\qquad \beta=L^2\gamma^{-1}\omega.
\end{equation}

Define $[\beta]_\times v=\beta\times v$. The normalization of the spatial structure constants gives

\begin{equation}\tag{M5}\label{eq:M5}
\begin{aligned}
\Omega_\beta&=4[\beta]_\times,&
\ell=\mathcal L_\beta\gamma&=\gamma\Omega_\beta-\Omega_\beta\gamma,\\
X&=\gamma^{-1}\gamma',&
Y&=\gamma^{-1}\ell,\qquad \mathrm{tr}Y=0.
\end{aligned}
\end{equation}

A prime denotes $d/dy$. The Lie derivative $\ell$ is generally nonzero for anisotropic shape. After integrating the total derivatives around the circle, the exact action is

\begin{equation}\tag{M6}\label{eq:M6}
\begin{aligned}
\mathcal S_{\mathrm{joint}}
&=\int_0^1dy\,[U+\mathcal L_{\partial}],\\
\mathcal L_{\partial}
&=\frac{\pi^2J}{32\kappa^2\mathcal N^2}
\left[\mathrm{tr}(X^2)-(\mathrm{tr}X)^2-2\mathrm{tr}(XY)\right].
\end{aligned}
\end{equation}

Here $U$ is the static expression (M3) evaluated at the instantaneous coordinates. Its shift-dependent static contribution is already included, so no extra $Y^2$ term is added to (M6). The first two trace terms form the velocity quadratic form; the last is linear in velocities. Torsion derivatives and mixing rotations contribute only total derivatives, displayed in Appendix J. The constant complete NY coupling therefore does not enter the periodic bulk action. Formula (M6) holds at finite amplitude within this spatially homogeneous Euclidean family; it is not a fourth-order truncation.
